\hypertarget{ds1-intro}{}

本章围绕两件事展开：区分数据结构中的基本概念；根据程序的执行过程分析时间与辅助空间。先理解定义和条件，再用少量代码验证，避免只背复杂度结论。


\hypertarget{ds1-concepts}{}

\section{数据结构：对象、关系与操作}\label{ds1-concepts}

\hypertarget{ds1-terms}{}

\begin{definition}[数据元素与数据项]
数据是信息的载体。\textbf{数据元素}是数据的基本单位，一个元素可由若干\textbf{数据项}组成；数据项是构成数据元素的不可分割的最小单位。\textbf{数据对象}是性质相同的数据元素的集合。
\end{definition}

\hypertarget{ds1-record}{}

例如，学生记录是数据元素，学号、姓名、成绩是其中的数据项；一批同类学生记录组成数据对象。讨论数据结构时，还要说明这些元素之间的关系。


\hypertarget{ds1-three}{}

\begin{takeaway}[三要素]
\textbf{数据结构的三要素：逻辑结构、存储结构、数据的运算。}逻辑结构描述元素之间的关系；存储结构把这些关系表示在计算机中；运算规定可以对数据做什么。
\end{takeaway}

\hypertarget{ds1-value-relation}{}

\Sidebar{值与关系}{存储时不仅保存元素值，还要体现元素之间的关系。关系可以由连续位置隐含，也可以由指针显式记录。}

\hypertarget{ds1-logical}{}

\par\smallskip\noindent\begin{tabularx}{\linewidth}{@{}YYY@{}}\toprule
\rowcolor{SNblue} \textbf{逻辑结构} & \textbf{元素关系} & \textbf{典型对象} \\
集合 & 只同属一个集合 & 无先后关系的对象集合 \\
线性结构 & 一对一的先后关系 & 线性表、栈、队列 \\
树形结构 & 一对多的层次关系 & 树、二叉树 \\
图状结构 & 多对多的一般关系 & 图 \\\bottomrule\end{tabularx}\par\smallskip

\hypertarget{ds1-linear-note}{}

在线性结构中，除首元素外，每个元素有唯一直接前驱；除尾元素外，每个元素有唯一直接后继。一个元素也可以同时是首元素和尾元素。


\hypertarget{ds1-storage}{}

\par\smallskip\noindent\begin{tabularx}{\linewidth}{@{}YY@{}}\toprule
\rowcolor{SNblue} \textbf{存储方式} & \textbf{如何体现关系} \\
顺序存储 & 逻辑相邻的元素保存在相邻存储位置 \\
链式存储 & 用结点内的指针连接逻辑相关的结点 \\
索引存储 & 建立索引项，记录定位信息 \\
散列存储 & 根据关键字计算存储位置 \\\bottomrule\end{tabularx}\par\smallskip

\hypertarget{ds1-not-equal}{}

\begin{pitfall}[关系与实现]
\textbf{逻辑结构与存储结构不能混为一谈。}线性表可以顺序存储，也可以链式存储；采用链式存储，也不能据此判定对象一定是一棵树。
\end{pitfall}

\hypertarget{ds1-adt}{}

\textbf{抽象数据类型（ADT）}描述数据对象、数据关系和基本操作，不限定具体存储实现。以线性表为例，“取第$i$个元素”是运算定义；用数组直接访问或沿链表查找，是两种实现。


\hypertarget{ds1-algorithm}{}

\section{算法的特性与设计要求}\label{ds1-algorithm}

\hypertarget{ds1-algorithm-def}{}

\begin{definition}[算法]
算法是求解一类问题的有限操作序列。对合法输入，应在有限步骤内结束，每一步含义明确且能够执行。
\end{definition}

\hypertarget{ds1-five}{}

\par\smallskip\noindent\begin{tabularx}{\linewidth}{@{}YY@{}}\toprule
\rowcolor{SNblue} \textbf{特性} & \textbf{准确含义} \\
有穷性 & 对每个合法输入，执行有限步后结束 \\
确定性 & 每条指令含义明确，不产生歧义 \\
可行性 & 每一步可用已实现的基本操作有限次完成 \\
输入 & 有零个或多个输入 \\
输出 & 有一个或多个输出 \\\bottomrule\end{tabularx}\par\smallskip

\hypertarget{ds1-output}{}

\Sidebar{输出是什么}{输出可以是返回值、修改后的数据状态等，不限于屏幕打印。}

\hypertarget{ds1-quality}{}

\textbf{正确性、可读性、健壮性以及时间与空间效率}是算法设计的重要评价要求。可行性不等于运行得足够快；有穷性也不等于执行次数很少。评价算法不能只看程序在某台计算机上运行了多少秒。


\hypertarget{ds1-max-boundary}{}

\begin{example}[边界条件（补充例）]
扫描数组求最大值时，若用\InlineCode{max=A[0]}初始化，就必须满足$n\ge1$。允许空数组时，应先处理空输入，不能访问不存在的首元素。
\end{example}

\hypertarget{ds1-time}{}

\section{时间复杂度：先计数，再判断数量级}\label{ds1-time}

\hypertarget{ds1-cost}{}

以输入规模$n$为变量，选择比较、赋值或数组访问等基本操作，统计其执行次数，再考察增长数量级。本章采用定长基本操作耗时为常数的模型。


\hypertarget{ds1-asymptotic}{}

\begin{takeaway}[记号与化简]
$O(f(n))$表示渐近上界，$\Theta(f(n))$表示紧确数量级。408选择题通常用$O$表达所要求的紧数量级；例如$3n^2+7n+9=\Theta(n^2)$，答题时通常写$O(n^2)$。
\end{takeaway}

\hypertarget{ds1-order}{}

常见增长次序为：$1\prec\log n\prec n\prec n\log n\prec n^2\prec n^3\prec2^n\prec n!$。这里表示渐近增长快慢，不是某个固定$n$下的运行时间排序。


\hypertarget{ds1-rules}{}

\par\smallskip\noindent\begin{tabularx}{\linewidth}{@{}YY@{}}\toprule
\rowcolor{SNblue} \textbf{程序结构} & \textbf{分析方法} \\
顺序执行的两个片段 & 次数相加，最终保留增长最快的项 \\
条件分支 & 按具体输入分析；求最坏情况时取最坏分支 \\
固定范围的嵌套循环 & 在每轮内层次数相同的条件下相乘 \\
内层范围依赖外层 & 逐轮写出内层次数，再求和 \\\bottomrule\end{tabularx}\par\smallskip

\hypertarget{ds1-count-method}{}

\Sidebar{计数的通用方法}{设循环执行$k$轮，把关键变量表示成$k$的函数，再代入终止条件。}

\hypertarget{ds1-patterns}{}

\par\smallskip\noindent\begin{tabularx}{\linewidth}{@{}YYY@{}}\toprule
\rowcolor{SNblue} \textbf{变化或次数} & \textbf{由$k$得到的关系} & \textbf{数量级} \\
每轮加1，直到$n$ & $k\approx n$ & $O(n)$ \\
每轮乘2，直到$n$ & $2^k\approx n$ & $O(\log n)$ \\
累计$1+2+\cdots+k$ & $k(k+1)/2\approx n$ & $O(\sqrt n)$ \\
内层依次执行$1,2,\ldots,n$ & $\sum_{i=1}^n i=n(n+1)/2$ & $O(n^2)$ \\\bottomrule\end{tabularx}\par\smallskip

\hypertarget{ds1-ex2017}{}

\subsection{累计和控制终止条件}\label{ds1-ex2017}

\hypertarget{ds1-ex2017-stem}{}

【2017年统考】下列函数的时间复杂度是（　）。


\hypertarget{ds1-code2017}{}

\Needspace{96.0pt}
\Lead{题目代码}
\CodeFile{code/complexity_2017.cpp}

\hypertarget{ds1-ex2017-options}{}

A. $O(\log n)$　B. $O(n^{1/2})$　C. $O(n)$　D. $O(n\log n)$


\hypertarget{ds1-ex2017-answer}{}

\begin{solution}
执行$k$轮后，$i=k$，$sum=k(k+1)/2$。循环结束时，$k$是使$k(k+1)/2\ge n$成立的最小非负整数，故$k=\Theta(\sqrt n)$，选B。\textbf{判断次数要看每轮如何增长，不能把最终的$sum\approx n$当成执行了$n$次。}
\end{solution}

\hypertarget{ds1-ex2022}{}

\subsection{两层循环也可能是线性时间}\label{ds1-ex2022}

\hypertarget{ds1-ex2022-stem}{}

【2022年统考】下列程序段的时间复杂度是（　）。


\hypertarget{ds1-code2022}{}

\Needspace{86.4pt}
\Lead{题目程序段}
\CodeFile{code/complexity_2022.cpp}

\hypertarget{ds1-ex2022-options}{}

A. $O(\log n)$　B. $O(n)$　C. $O(n\log n)$　D. $O(n^2)$


\hypertarget{ds1-ex2022-answer}{}

\begin{solution}
当$n>1$时，各轮内层次数为$1,2,4,\ldots,2^k$，其中$2^k<n\le2^{k+1}$。总次数$1+2+\cdots+2^k=2^{k+1}-1=\Theta(n)$，选B。把每轮都按内层最大次数$n$估计，会得到不紧的上界。
\end{solution}

\hypertarget{ds1-cases}{}

\subsection{最好、最坏与平均情况}\label{ds1-cases}

\hypertarget{ds1-cases-text}{}

同样长度为$n$的无序表，顺序查找在首位命中时只比较1次，最坏需要比较$n$次。最好和最坏分别取同规模合法输入中的最小、最大代价；平均情况还必须说明概率分布。


\hypertarget{ds1-average}{}

若保证查找成功，且每个元素被查找的概率相同，平均比较次数为$(1+2+\cdots+n)/n=(n+1)/2$，仍是$O(n)$。平均代价一般不能直接取最好与最坏的算术平均。


\hypertarget{ds1-space}{}

\section{辅助空间：看同时保留多少状态}\label{ds1-space}

\hypertarget{ds1-space-def}{}

\begin{definition}[辅助空间]
辅助空间复杂度衡量算法额外使用的空间随输入规模增长的情况，通常不计输入本身。输出区域是否计入，按题目约定处理。
\end{definition}

\hypertarget{ds1-space-table}{}

\par\smallskip\noindent\begin{tabularx}{\linewidth}{@{}YYY@{}}\toprule
\rowcolor{SNblue} \textbf{场景} & \textbf{辅助空间} & \textbf{原因} \\
少量循环变量 & $O(1)$ & 数量不随$n$增长 \\
额外申请$n$个单元 & $O(n)$ & 需同时保留$n$个单元 \\
递归深度$n$，每层常数空间 & $O(n)$ & 调用栈尚未返回 \\
非递归函数依次调用$n$次 & $O(1)$* & 前次返回后可复用 \\\bottomrule\end{tabularx}\par\smallskip

\hypertarget{ds1-space-star}{}

*最后一行假设每次调用只用常数局部空间，且没有把结果累积保存到增长的容器中。


\hypertarget{ds1-constant}{}

\Sidebar{常数空间}{$O(1)$不等于不用空间，而是额外空间不随输入规模增长。}

\hypertarget{ds1-recursion}{}

例如递归求阶乘，每次调用只做常数工作，但调用尚未返回时，参数、返回位置等状态要保留。按题目常用的单位成本模型，时间为$O(n)$，递归栈空间也为$O(n)$。\textbf{总调用次数影响时间，最大同时调用深度影响栈空间。}


\hypertarget{ds1-inplace}{}

原地算法通常只允许常数辅助空间。数组在原位置修改，并不自动意味着空间是$O(1)$；若使用递归，仍要计入调用栈。


\hypertarget{ds1-recall}{}

\begin{recallcard}[检查是否理解条件]
\textbf{基础自检（补充题）}

1. 线性表是否必须连续存储？\par
2. 每轮把变量乘2，直到达到$n$，需要多少轮？\par
3. 循环执行$n$次，是否就一定使用$O(n)$辅助空间？

\textbf{答案：}1. 否，线性表描述逻辑关系，也可以链式存储。2. $O(\log n)$轮。3. 否，空间看同时占用的额外存储，几个循环变量仍是$O(1)$。
\end{recallcard}

\chapter{线性表}
\hypertarget{ds2-intro}{}

本章重点是线性表的基本概念、两种存储实现和基本操作。所有代码采用C++17写法：位序从1开始，数组下标从0开始；普通单链表带有效头结点，尾结点的next为空。


\hypertarget{ds2-basics}{}

\section{定义、位序与基本操作}\label{ds2-basics}

\hypertarget{ds2-definition}{}

\begin{definition}[线性表]
线性表是具有相同数据类型的$n$个数据元素的有限序列。$n=0$时为空表；非空表的首元素没有直接前驱，尾元素没有直接后继，其余元素各有唯一直接前驱和直接后继。
\end{definition}

\hypertarget{ds2-values}{}

“相同数据类型”不要求数值互不相同；“有限序列”同时强调元素数量有限和次序确定。只有一个元素时，该元素既没有直接前驱，也没有直接后继。


\hypertarget{ds2-operations}{}

\par\smallskip\noindent\begin{tabularx}{\linewidth}{@{}YY@{}}\toprule
\rowcolor{SNblue} \textbf{基本操作} & \textbf{含义或条件} \\
初始化、判空、求长 & 建立空表；判断$n=0$；取得元素个数 \\
按位序取值 & 取第$i$个元素，$1\le i\le n$ \\
按值查找 & 查找满足给定值条件的元素 \\
插入 & 在第$i$个位置插入，$1\le i\le n+1$ \\
删除 & 删除第$i$个元素，$1\le i\le n$ \\
销毁 & 释放动态申请的存储空间 \\\bottomrule\end{tabularx}\par\smallskip

\hypertarget{ds2-index}{}

\begin{takeaway}[位序与下标]
\textbf{第$i$个元素对应数组下标$i-1$。}长度为$n$的表有$n+1$个合法插入位置，却只有$n$个合法删除位置。插入第$n+1$位表示追加到表尾。
\end{takeaway}

\hypertarget{ds2-find-failure}{}

\Sidebar{失败标记}{按位序返回时，可用0表示查找失败；若按数组下标返回，应使用-1等独立标记，避免与下标0混淆。}

\hypertarget{ds2-array}{}

\section{顺序表：连续存储与元素移动}\label{ds2-array}

\hypertarget{ds2-seq-def}{}

顺序表用一段连续存储空间依次保存元素。若首元素地址为$b$，每个元素占$w$个地址单位，第$i$个元素地址为$b+(i-1)w$，因此可以在$O(1)$时间按位序访问。


\hypertarget{ds2-random}{}

\begin{pitfall}[按位序与按值]
\textbf{随机存取指按给定位序直接定位。}无序顺序表按值查找，仍然要比较元素，最坏时间为$O(n)$。
\end{pitfall}

\hypertarget{ds2-sq-struct}{}

\Needspace{96.0pt}
\Lead{顺序表的静态表示}
\CodeFile{code/sq_struct.cpp}

\hypertarget{ds2-capacity}{}

\InlineCode{length}是有效元素个数，\InlineCode{MaxSize}是容量，始终满足$0\le length\le MaxSize$。静态数组容量固定，表长可以变化；动态分配的数组仍需保持有效元素连续，扩容时可能搬移原有元素。


\hypertarget{ds2-sq-find}{}

\Needspace{96.0pt}
\Lead{按值查找：返回第一个匹配元素的位序，失败返回0}
\CodeFile{code/sq_find.cpp}

\hypertarget{ds2-insert}{}

\subsection{插入：先从后向前移动}\label{ds2-insert}

\hypertarget{ds2-insert-idea}{}

在第$i$位插入$x$，要先将原第$i$到第$n$个元素各后移一位，空出\InlineCode{data[i-1]}。移动必须从表尾开始，避免覆盖尚未搬走的元素。


\hypertarget{ds2-array-figure}{}

\MainFigure{在第2位插入99：先把30移到下标3，再把20移到下标2，最后写入99。}{
\begin{tikzpicture}[x=11mm,y=11mm,font=\small,
  cell/.style={draw=SNink,minimum width=11mm,minimum height=8mm,inner sep=0pt},
  lab/.style={font=\scriptsize,text=SNmuted},>=Latex]
\node[anchor=east] at (-.7,0) {before};
\node[anchor=east] at (-.7,-1.9) {after};
\foreach \v [count=\i from 0] in {10,20,30,{}} {
  \node[cell,fill=SNblue] (b\i) at (\i,0) {\v};
  \node[lab] at (\i,.65) {\i};}
\foreach \v [count=\i from 0] in {10,99,20,30} {
  \node[cell,fill=SNblue] (a\i) at (\i,-1.9) {\v};}
\node[cell,fill=SNpink] at (1,-1.9) {99};
\draw[->,thick,SNrose] (b2.south) -- (a3.north);
\draw[->,thick,SNink] (b1.south) -- (a2.north);
\node[lab,anchor=west] at (3.65,0) {$n=3$};
\node[lab,anchor=west] at (3.65,-1.9) {$n=4$};
\end{tikzpicture}
}

\hypertarget{ds2-sq-insert}{}

\Needspace{134.4pt}
\Lead{按位序插入：空间不足或位序非法时不修改表}
\CodeFile{code/sq_insert.cpp}

\hypertarget{ds2-insert-boundary}{}

\Sidebar{插入边界}{空表只能在第1位插入。尾插不移动元素；头插移动全部$n$个元素。这里分析的数组已有可用容量。}

\hypertarget{ds2-insert-count}{}

移动元素数为$n-i+1$。最好情况为尾插，时间$O(1)$；最坏为头插，时间$O(n)$。若$n+1$个位置等概率选择，平均移动次数为$\dfrac{n+(n-1)+\cdots+0}{n+1}=\dfrac n2$，平均时间为$O(n)$。


\hypertarget{ds2-delete}{}

\subsection{删除：保存元素，再从前向后移动}\label{ds2-delete}

\hypertarget{ds2-delete-idea}{}

删除第$i$个元素时，先保存\InlineCode{data[i-1]}，再把原第$i+1$至第$n$个元素依次前移，最后表长减1。有效区之外的旧数值不影响表的逻辑内容。


\hypertarget{ds2-sq-delete}{}

\Needspace{124.8pt}
\Lead{按位序删除：x返回被删除的值}
\CodeFile{code/sq_delete.cpp}

\hypertarget{ds2-delete-count}{}

移动元素数为$n-i$。尾删为$O(1)$，头删最坏为$O(n)$。若$n$个已有元素等概率被删除，平均移动次数为$\dfrac{(n-1)+(n-2)+\cdots+0}{n}=\dfrac{n-1}{2}$。


\hypertarget{ds2-moves-example}{}

\begin{example}[移动次数（补充例）]
长度为8的顺序表，在第3位插入需移动$8-3+1=6$个元素；删除原表第3个元素需移动$8-3=5$个元素。两种操作的合法位置数不同，平均移动次数的分母也不同。
\end{example}

\hypertarget{ds2-linked}{}

\section{单链表：结点、查找与链接修改}\label{ds2-linked}

\hypertarget{ds2-node-def}{}

单链表的结点包含数据域和后继指针域。不同结点之间不要求地址连续，逻辑次序由\InlineCode{next}决定；同一个结点作为整体分配，其内部存储单元是连续的。


\hypertarget{ds2-node-struct}{}

\Needspace{105.6pt}
\Lead{结点类型与带头结点的空表}
\CodeFile{code/node_struct.cpp}

\hypertarget{ds2-head-figure}{}

\MainFigure{L是头指针，指向头结点；首元结点为L->next。头结点不保存有效表元素。}{
\begin{tikzpicture}[x=1mm,y=1mm,font=\small,>=Latex,
 nodebox/.style={draw=SNink,fill=SNblue,minimum width=14mm,minimum height=9mm,inner sep=1mm}]
\node (L) at (0,0) {$L$};
\node[nodebox] (h) at (19,0) {head};
\node[nodebox] (a) at (45,0) {10};
\node[nodebox] (b) at (70,0) {20};
\node (null) at (96,0) {$\varnothing$};
\draw[->] (L)--(h);\draw[->] (h)--(a);\draw[->] (a)--(b);\draw[->] (b)--(null);
\node[font=\scriptsize,text=SNmuted] at (45,-10) {first data node};
\end{tikzpicture}
}

\hypertarget{ds2-head-terms}{}

\par\smallskip\noindent\begin{tabularx}{\linewidth}{@{}YY@{}}\toprule
\rowcolor{SNblue} \textbf{名称} & \textbf{指的是什么} \\
头指针 & 保存链表起始结点地址的指针变量 \\
头结点 & 为统一操作而附加的结点，不计入表长 \\
首元结点 & 保存第一个有效数据元素的结点 \\\bottomrule\end{tabularx}\par\smallskip

\hypertarget{ds2-head-benefit}{}

带头结点后，首位置的插入和删除也有统一的前驱结点。普通带头结点单链表为空时，\InlineCode{L}仍指向有效头结点，满足\InlineCode{L->next == nullptr}。无头结点的空表则通常用头指针为空表示。


\hypertarget{ds2-linked-find}{}

\subsection{按位序查找与按值查找}\label{ds2-linked-find}

\hypertarget{ds2-getelem}{}

\Needspace{115.2pt}
\Lead{本函数约定i=0返回头结点，i>=1返回第i个数据结点}
\CodeFile{code/get_elem.cpp}

\hypertarget{ds2-getelem-note}{}

合法数据位序为$1\sim n$；越过表尾返回空指针。开放$i=0$是为了方便寻找插入、删除位置的前驱，不表示头结点属于数据元素。按位序查找最坏时间为$O(n)$。


\hypertarget{ds2-linked-locate}{}

\Needspace{105.6pt}
\Lead{按值查找：返回第一个匹配结点，失败返回空指针}
\CodeFile{code/linked_locate.cpp}

\hypertarget{ds2-linked-insert}{}

\subsection{插入：先保存旧后继，再接入新结点}\label{ds2-linked-insert}

\hypertarget{ds2-linked-insert-idea}{}

在第$i$位插入前，先找到第$i-1$个结点\InlineCode{p}。新结点先接到旧后继，再让\InlineCode{p}接到新结点；两句赋值不能直接颠倒。


\hypertarget{ds2-link-insert-figure}{}

\MainFigure{插入分两步：s->next接旧后继q；p->next再接s。若先覆盖p->next，可能丢失原来的后半链。}{
\begin{tikzpicture}[x=1mm,y=1mm,font=\small,>=Latex,
 nodebox/.style={draw=SNink,fill=SNblue,minimum width=13mm,minimum height=8mm,inner sep=1mm}]
\node[nodebox] (p) at (12,0) {$p$};
\node[nodebox] (q) at (76,0) {$q$};
\node[nodebox,fill=SNpink] (s) at (44,-17) {$s$};
\draw[->,dashed,SNmuted] (p.east)--node[above,font=\scriptsize]{old next}(q.west);
\draw[->,thick,SNink] (s.east)--node[below right,font=\scriptsize]{1: s->next}(q.south west);
\draw[->,thick,SNrose] (p.south east)--node[below left,font=\scriptsize]{2: p->next}(s.west);
\end{tikzpicture}
}

\hypertarget{ds2-linked-insert-code}{}

\Needspace{124.8pt}
\Lead{按位序插入：i=1时前驱就是头结点}
\CodeFile{code/linked_insert.cpp}

\hypertarget{ds2-linked-allocation}{}

\Sidebar{C/C++表示}{代码使用new申请结点、delete释放结点；C语言可用malloc/free实现同一结构，分配与释放方式要配对。}

\hypertarget{ds2-linked-insert-cost}{}

已知有效前驱时，修改链接只需$O(1)$；从头寻找第$i-1$个结点则最坏需要$O(n)$。因此按位序插入的完整操作最坏仍为$O(n)$。


\hypertarget{ds2-linked-delete}{}

\subsection{删除：先接好后链，再释放结点}\label{ds2-linked-delete}

\hypertarget{ds2-linked-delete-code}{}

\Needspace{153.6pt}
\Lead{按位序删除：先检查前驱和待删结点是否存在}
\CodeFile{code/linked_delete.cpp}

\hypertarget{ds2-delete-pitfall}{}

\begin{pitfall}[删除的前提]
释放\InlineCode{q}之后，不能再读取\InlineCode{q->next}。只有待删结点的指针、却没有前驱时，普通单链表通常仍需从头找前驱；复制后继值再删后继的技巧不适用于尾结点，也不一定符合题目对结点身份的要求。
\end{pitfall}

\hypertarget{ds2-build}{}

\subsection{建立链表：头插与尾插}\label{ds2-build}

\hypertarget{ds2-head-build}{}

\Needspace{105.6pt}
\Lead{头插法：调用前L是空表，按输入逆序建立}
\CodeFile{code/head_build.cpp}

\hypertarget{ds2-tail-build}{}

\Needspace{124.8pt}
\Lead{尾插法：调用前L是空表，按输入原序建立}
\CodeFile{code/tail_build.cpp}

\hypertarget{ds2-build-order}{}

输入$1,2,3$，头插得到$3\to2\to1$，尾插得到$1\to2\to3$。两者均只对每个元素做常数次链接操作，建立时间为$O(n)$。尾插若每次都从头找尾部，才会退化为$O(n^2)$。


\hypertarget{ds2-build-space}{}

建立$n$个数据结点本身需要$O(n)$存储；除结果链表外，工作指针只占$O(1)$辅助空间。回答空间复杂度时，要说明是否把结果结构计入。


\hypertarget{ds2-destroy}{}

\Needspace{115.2pt}
\Lead{销毁普通无环单链表：包含头结点，结束后L为空}
\CodeFile{code/destroy.cpp}

\hypertarget{ds2-variants}{}

\section{双链表、循环链表与静态链表}\label{ds2-variants}

\hypertarget{ds2-double}{}

\subsection{双链表：同时维护前后链接}\label{ds2-double}

\hypertarget{ds2-double-text}{}

双链表的结点增加前驱指针\InlineCode{prior}。已知结点后，可以在$O(1)$时间访问它的直接前驱；按位序定位结点仍通常需要$O(n)$时间。


\hypertarget{ds2-double-insert}{}

\Needspace{153.6pt}
\Lead{在p之前插入s：先设置新结点两侧，再修改旧结点链接}
\CodeFile{code/double_insert.cpp}

\hypertarget{ds2-double-invariant}{}

设旧前驱为$r$，插入后必须同时满足$r.next=s$、$s.prior=r$、$s.next=p$、$p.prior=s$。代码假设\InlineCode{p->prior}有效；无头结点链表在首元结点之前插入时，应另行修改头指针。


\hypertarget{ds2-double-delete}{}

删除具有有效前驱和后继的数据结点\InlineCode{p}时，先接好两侧，再释放\InlineCode{p}。下面的操作假设两个邻接结点都有效；普通非循环链表的尾结点后继为空，需另作判空处理。


\hypertarget{ds2-double-delete-code}{}

\Needspace{76.8pt}
\Lead{双链表删除的链接操作}
\CodeFile{code/double_delete.cpp}

\hypertarget{ds2-circular}{}

\subsection{循环链表：终止位置变成头结点}\label{ds2-circular}

\hypertarget{ds2-circular-text}{}

带头结点的循环单链表，尾结点的\InlineCode{next}指回头结点。空表时头结点指向自己；遍历遇到头结点时结束，而不是等待空指针。


\hypertarget{ds2-circular-figure}{}

\MainFigure{带头结点的普通单链表与循环单链表：空表链接不同，遍历终止条件也不同。}{
\begin{tikzpicture}[x=1mm,y=1mm,font=\small,>=Latex,
 nodebox/.style={draw=SNink,fill=SNblue,minimum width=15mm,minimum height=8mm,inner sep=1mm}]
\node[anchor=east,font=\scriptsize] at (23,0) {ordinary};
\node (l1) at (32,0) {$L$}; \node[nodebox] (h1) at (54,0) {head};
\node (null) at (86,0) {$\varnothing$};
\draw[->] (l1)--(h1);\draw[->] (h1)--(null);
\node[anchor=east,font=\scriptsize] at (23,-22) {circular};
\node (l2) at (32,-22) {$L$};\node[nodebox] (h2) at (54,-22) {head};
\draw[->] (l2)--(h2);
\draw[->,SNrose,thick] (h2.east) -- ++(15,0) -- ++(0,-11) -| (h2.south);
\end{tikzpicture}
}

\hypertarget{ds2-empty}{}

\par\smallskip\noindent\begin{tabularx}{\linewidth}{@{}YY@{}}\toprule
\rowcolor{SNblue} \textbf{结构约定} & \textbf{空表条件} \\
普通单链表，带有效头结点 & \InlineCode{L->next == nullptr} \\
循环单链表，带有效头结点 & \InlineCode{L->next == L} \\
循环双链表，带有效头结点 & \InlineCode{L->next == L}且\InlineCode{L->prior == L} \\\bottomrule\end{tabularx}\par\smallskip

\hypertarget{ds2-pointer-address}{}

\Sidebar{L与\&L}{\InlineCode{L}是指针保存的结点地址；\InlineCode{&L}是指针变量自身的地址。二者不能替换。}

\hypertarget{ds2-tail-circle}{}

循环单链表保存尾指针\InlineCode{rear}时，可通过\InlineCode{rear->next}取得链头。对两个非空、无头结点的循环单链表，交换尾部的后继链接即可$O(1)$拼接；只有头指针时，通常还要线性查找尾结点。带头结点的拼接还须处理多余头结点。


\hypertarget{ds2-static}{}

\subsection{静态链表：用数组下标代替指针}\label{ds2-static}

\hypertarget{ds2-static-text}{}

静态链表把结点存放在数组中，用整数游标记录后继结点的数组下标。数组中相邻位置不一定是逻辑后继；它仍按链接查找，不因此获得按位序的常数时间随机访问。通常还需管理空闲结点链。


\hypertarget{ds2-compare}{}

\par\smallskip\noindent\begin{tabularx}{\linewidth}{@{}YYY@{}}\toprule
\rowcolor{SNblue} \textbf{维度} & \textbf{顺序表} & \textbf{普通单链表} \\
地址要求 & 元素连续 & 不同结点可分散 \\
关系表示 & 由位置隐含 & 由next显式记录 \\
按位序访问，最坏 & $O(1)$ & $O(n)$ \\
无序按值查找，最坏 & $O(n)$ & $O(n)$ \\
任意位序插删，最坏 & $O(n)$，需移动 & $O(n)$，含定位 \\
额外结构空间 & 无需链接域 & 每个结点有指针域 \\\bottomrule\end{tabularx}\par\smallskip

\hypertarget{ds2-choose}{}

\begin{takeaway}[选择结构的依据]
频繁按位序访问，优先考虑顺序表。若经常在已定位的位置修改链接，链式结构较方便。选择结构时，应同时考虑定位成本、修改成本和存储开销。
\end{takeaway}

\hypertarget{ds2-basic-algorithms}{}

\section{基础巩固：一次扫描与原地逆置}\label{ds2-basic-algorithms}

\hypertarget{ds2-delete-all}{}

\subsection{顺序表删除所有值为x的元素}\label{ds2-delete-all}

\hypertarget{ds2-delete-all-idea}{}

使用读下标$j$扫描原表，用$k$记录已保留元素个数。遇到不等于$x$的元素，就写到保留区末尾。这样不需要每删一个元素就搬移后半表。


\hypertarget{ds2-delete-all-code}{}

\Needspace{115.2pt}
\Lead{保持原有相对次序，时间O(n)，辅助空间O(1)}
\CodeFile{code/delete_all.cpp}

\hypertarget{ds2-delete-all-boundary}{}

整个过程中$k\le j$，写入不会破坏尚未读取的元素；全部删除时$k=0$，全部保留时$k=n$。


\hypertarget{ds2-reverse}{}

\subsection{带头结点单链表原地逆置}\label{ds2-reverse}

\hypertarget{ds2-reverse-idea}{}

\InlineCode{prev}指向已逆置部分，\InlineCode{cur}指向未处理部分。每轮先保存旧后继，再反转当前链接，随后推进工作指针。


\hypertarget{ds2-reverse-code}{}

\Needspace{144.0pt}
\Lead{保留头结点，时间O(n)，辅助空间O(1)}
\CodeFile{code/reverse.cpp}

\hypertarget{ds2-reverse-boundary}{}

空表时循环不执行，\InlineCode{L->next}仍为空；单结点表也能正常处理。最后要更新头结点的\InlineCode{next}，不能把头结点当作数据结点一起逆置。


\hypertarget{ds2-recall}{}

\begin{recallcard}[用条件检验结论]
\textbf{基础自检（补充题）}

1. 顺序表随机存取是否表示按值查找一定为$O(1)$？\par
2. 长度为8的顺序表，插入第3位、删除第3个元素，分别移动多少个元素？\par
3. 已知单链表前驱后删除其后继，修改链接要多久？若先按位序定位呢？\par
4. 带头结点循环单链表为空时，\InlineCode{L->next}等于什么？\par
5. 头插法和维护尾指针的尾插法，建立$n$个结点各需要多少时间？

\textbf{答案：}1. 否，无序表按值查找最坏$O(n)$。2. 分别为6和5。3. 修改链接$O(1)$；含定位最坏$O(n)$。4. 等于\InlineCode{L}。5. 都是$O(n)$；输入次序分别逆转和保持。
\end{recallcard}
